\subsection{The complete colimit criterion used in C.28}
\begin{theorem}[Step audit of INDLOCAL C.27 and PRETOT F.1/F.4]
\label{F:C27-colimits}
For a small index category $A$, the canonical functor
\[
 \int_A D\longrightarrow\colim_A^{\Cat}D
\]
is a nerve weak equivalence for every functor $D:A\to\Cat$
if and only if every connected component of $A$ is filtered.
The empty index is included. In particular the sequential indices
in the cited manuscript passages satisfy the criterion, whereas
the opposite type category does not.
\end{theorem}
\begin{proof}
\textbf{The filtered implication and reduction to one component.}
For filtered $A$, finite strings, their composition equations,
and finite simplicial homotopy data occur at one stage. Thus
nerve commutes with the colimit and the homotopy-colimit comparison
is a weak equivalence, as proved in Proposition~\ref{F:C22-types}.
Thomason's Grothendieck comparison identifies its source with
$B\int_A D$. All the constructions split as disjoint unions over
components, including the empty case. For necessity it is therefore
enough to take $A$ connected and nonempty. The constant one-point
diagram gives $BA\simeq *$.

\textbf{Common successors, checking the set-valued detection.}
For a discrete-valued diagram the conclusion says that each
component of its category of elements is weakly contractible.
If $a,b$ have no common successor, let $U_a,U_b$ be their principal
cosieves. They are disjoint nonempty connected full subcategories:
every object in each is reached from its named object. Form the
two-point diagram which identifies its two points exactly on
$U_a\amalg U_b$. Its element category is two copies of $A$ glued
along those two cosieves. This also describes the nerves as a union
of subcomplexes, because after a nerve string enters a cosieve it
cannot leave. The two copies have contractible realization. The
kernel of
$H_0(BU_a\amalg BU_b)\to H_0(BA)\oplus H_0(BA)$ is a copy of
$\mathbb Z$, so Mayer--Vietoris gives nonzero first homology in
the connected union. Its ordinary colimit is a singleton, a
contradiction. Every pair of objects has a common successor.

\textbf{The group-valued diagram for a nonequalizable pair.}
Suppose $u,v:p\rightrightarrows q$ have no post-equalizer. Write
$P=A(p,-)$ and $Q=A(q,-)$. At $x$, make a bipartite graph whose
two vertex colors are $Q(x)$ and $P(x)$ and whose edges join
$f\in Q(x)$ to $fu$ and to $fv$. These two edges are distinct:
an equality would itself be a post-equalizer. Let $G(x)$ have
the group presentation
\[
 \left\langle Q(x)\amalg P(x)\ \middle|\
       [f,fu]=[f,fv]=1\text{ for }f\in Q(x)\right\rangle.
\]
Postcomposition preserves the colors and these relations, so this
is an actual group-valued functor on $A$.

\textbf{The primary theorem and every required hypothesis.}
For a finite graph the Salvetti complex is aspherical by
\cite[Theorem~2.6 and Section~2.7]{charney2007}; its signed vertex
links are flag, and its universal cover is contractible. Here the
graph is bipartite, so its only positive-dimensional cells are
the generator circles and one torus cell for each graph edge.
Their cellular boundaries are zero. The same conclusions for an
infinite graph follow by the union of its finite full subgraphs:
each sphere and nullhomotopy uses finitely many cells. Orient every
torus by the $Q$ coordinate followed by the $P$ coordinate.
Transitions preserve these colors, so they send each edge class
to its image edge with coefficient $+1$, even when distinct
vertices of the same color are identified. Therefore, as actual
homology diagrams,
\[
 H_0(BG(-))=\mathbb Z,\qquad
 H_1(BG(-))=\mathbb Z[P\amalg Q],\qquad
 H_2(BG(-))=\mathbb Z[Q\amalg Q],
\]
with all higher groups zero. This checks the naturality needed
for the argument; a list of objectwise homology groups alone
would not suffice.

\textbf{The two specified second homology groups disagree.}
Free abelian representables are projective in the category of
covariant abelian-group diagrams and have colimit $\mathbb Z$.
The derived colimits of the constant diagram are the homology of
$BA$, which is contractible. Thus in the bar double-complex
spectral sequence for $\hocolim_A BG(x)$ every positive
derived-colimit column vanishes. The displayed degree-two row
gives
\[
 H_2\bigl(\hocolim_A BG(x);\mathbb Z\bigr)=\mathbb Z^2.
\]
The ordinary categorical colimit of the one-object groups is
their one-object colimit group: connectedness identifies the
objects and each inverse equation survives the colimit
presentation. Both representable sets have singleton colimit,
so that group has exactly two generators and their commutator
relation. It is $\mathbb Z^2$, whose classifying space has second
homology $\mathbb Z$. The canonical comparison cannot be a weak
equivalence. This contradiction proves post-equalization of every
parallel pair, and hence filteredness.

\textbf{The inventory in the manuscript.}
The sequences have the actual index $\mathbb N$; using a fibred
description of a total category does not reverse that colimit
index. The opposite type category fails filteredness by the two
outer-position maps checked in Proposition~\ref{F:C22-types}.
The stated interval pushout has the direct free-category
universal property verified there, and its two intervals, common
endpoint and concatenation are contractible. This proves that
particular homotopy-pushout comparison without asserting the
universal property for every diagram on a span. These observations
check every use of C.27 in C.28.
\end{proof}
